\VerdierSemanticAnchor{front-title}
\VerdierPhysicalPageAnchor{0002}
\begin{titlepage}
  \centering
  {\Huge\bfseries 239\hfill ASTÉRISQUE\par}
  \vspace{1.5em}
  {\LARGE\bfseries 1996\par}
  \vfill
  {\Huge\bfseries 阿贝尔范畴的\par}
  \vspace{0.4em}
  {\Huge\bfseries 导出范畴\par}
  \vspace{3em}
  {\Large\bfseries Jean-Louis VERDIER\par}
  \vfill
  {\Large\bfseries 法国数学会\par}
  \smallskip
  {\small\bfseries 本书由法国国家科学研究中心资助出版\par}
\end{titlepage}

\VerdierSemanticAnchor{front-ams-classification}
\VerdierPhysicalPageAnchor{0003}
\chapter*{AMS 主题分类}
\addcontentsline{toc}{chapter}{AMS 主题分类}

14F20, 18E30, 18E35, 18G05, 18G10, 18G15, 18G20, 18G35,
18G40, 55U15, 55U25, 55U30.

\clearpage
\VerdierSemanticAnchor{front-portrait}
\VerdierPhysicalPageAnchor{0004}
\thispagestyle{empty}
\vspace*{0.62\textheight}
\begin{center}
Jean-Louis VERDIER，1977
\end{center}

\clearpage
\VerdierPhysicalPageAnchor{0005}
\thispagestyle{empty}
\null
\clearpage

\VerdierSemanticAnchor{front-preface}
\VerdierPhysicalPageAnchor{0006}
\chapter*{序言}
\addcontentsline{toc}{chapter}{序言}

Jean-Louis Verdier 于 1989 年 8 月 25 日不幸逝世。他曾于 1967 年 6 月
14 日在巴黎答辩其国家博士论文，题为《阿贝尔范畴的导出范畴》。本书采用的
文本就是该论文打字稿的最终版本。编者所做的工作仅限于改正细小错误。提交给
学院的只有题为“引言”的部分；正文在答辩时尚未完成，此后也一直如此。

这项工作的由来已相当为人所知（\emph{参见} [I]）。我们知道，Grothendieck
在二十世纪六十年代初提出了导出范畴的思想，目的是为同调代数提供一个框架，
以实现他在 1958 年爱丁堡国际数学家大会报告之后所设想的对偶定理之宏大推广。
他把建立这一形式体系作为 Verdier 的博士论文课题。1963 年，Verdier 发表了
一册概述该理论核心内容的结果汇编《导出范畴：若干结果（状态 0）》，由 IHÉS
油印（后收入 [SGA 4 1/2]）。导出范畴的概念固然源自 Grothendieck，但更早
提出三角范畴概念的则是 Verdier。那里给出的公理体系受稳定同伦论启发
（\emph{参见} [P]），后来证明具有出人意料的丰硕成果。

又过了若干年，Verdier 才着手全面整理这项工作。他为什么没有完成？人们提出
过许多猜想。我倾向于以下解释：缺失的部分正是导出函子理论。当 Verdier 写到
这里时，导出范畴已经在代数几何中得到广泛使用，各种变体和推广也在发展：
Deligne 式的逐点可导性、非加性的导出函子（Quillen 式同伦代数），以及滤过
导出函子。此时，那册结果汇编所确定的框架已显得不足。要找到恰当的观点，必须
重新开展一项基础性工作；而当时正关注其他问题的 Verdier，在我看来既无暇，
大概也无意去完成它。随着时间推移，关于导出范畴的论文不断出现，其应用日益
普遍，这一形式体系还进入了数学的其他领域（例如线性偏微分方程）。于是，在
Verdier 眼中，完成这部著作越来越成为一项没有必要的工作……

\clearpage
\VerdierPhysicalPageAnchor{0007}
这部文本虽不完整，却仍是一部珍贵乃至无可替代的参考资料（例如，其中对谱序列
所作的原创而优雅的论述，据我所知，在别处无从见到）。它终于得以问世，实在
令人欣慰。

\begin{flushright}
巴黎，1996 年 4 月 27 日\\
Luc Illusie
\end{flushright}

\section*{参考文献}

\begin{description}[labelwidth=6.2em,leftmargin=6.8em]
\item[{[I]}] L. Illusie，\emph{Catégories dérivées et dualité, travaux de
J.-L. Verdier}，L'Enseignement Mathématique，36 (1990)，369--391。
（1989 年 10 月 19 日，Université Paris VII 举办 Jean-Louis Verdier
纪念仪式时所作的报告。）

\item[{[P]}] D. Puppe，\emph{On the formal structure of stable homotopy
theory}，Colloquium on Algebraic Topology，Aarhus Universitet (1962)，
65--71。

\item[{\mbox{[SGA 4 1/2]}}] “Cohomologie étale”，Séminaire de Géométrie
Algébrique du Bois-Marie，SGA 4 1/2，P. Deligne 著，Lecture Notes in
Mathematics 569，Springer-Verlag (1977)，262--311。
\end{description}
